% source: RH-Framework/theorems/T01_NATIVE_SUMMABILITY_AND_RECOGNITION_COMPLETION.md
\hypertarget{t01-abc-native-cut-tail-and-recognition-cauchy-completion}{%
\chapter{T01-A/B/C --- Native Cut-Tail and Recognition-Cauchy Completion}\label{t01-abc-native-cut-tail-and-recognition-cauchy-completion}}

\textbf{Classification:} \texttt{PROVED}

\textbf{Foundational origin:} \texttt{NATIVE\_DERIVED}

\textbf{Main T01 classification:} \texttt{OPEN}

This theorem removes the first functional-analysis import from the recognition
foundation. It does not begin with

\begin{verbatim}
ell1;
ell2;
Hilbert space;
operator norm;
Haar measure;
Minkowski, Young, or an imported Cauchy-Schwarz theorem.
\end{verbatim}

It begins with finite typed residue paths over exact oriented-cut scalars and
defines two independent native completion relations.

The meta-foundation remains ordinary finite algebra, rational order, sequences,
and equivalence relations. No claim of foundation-free mathematics is made.
The relevant claim is narrower and precise: no pre-existing complex Hilbert or
Banach structure generates the objects proved below.

\begin{center}\rule{0.5\linewidth}{0.5pt}\end{center}

\hypertarget{primitive-finite-path-core}{%
\section{1. Primitive finite path core}\label{primitive-finite-path-core}}

Let

\[
\mathbb K_\Sigma^0
=
\{r+t\iota_\Sigma:r,t\in\mathbb Q_\Sigma\}
\]

be the exact cut field, with

\[
\iota_\Sigma^2=-1,
\qquad
(r+t\iota_\Sigma)^\dagger=r-t\iota_\Sigma.
\]

Define the native radial square

\[
N_\Sigma(z)=z^\dagger z=r^2+t^2
\in\mathbb Q_{\Sigma,+}.
\tag{1.1}
\]

Let \(\mathscr G_\Sigma\) be a typed residue groupoid. Its object and arrow sets
may be countable, but every element used in the dense core has finite support.
Write

\[
\mathscr P_\Sigma^0
=
C_c(\mathscr G_\Sigma,\mathbb K_\Sigma^0).
\]

Convolution and dagger are defined entirely by path composition and path
reversal:

\[
(a\star b)_\gamma
=
\sum_{\alpha\beta=\gamma}a_\alpha b_\beta,
\tag{1.2}
\]

\[
(a^\dagger)_\gamma
=
(a_{\gamma^{-1}})^\dagger.
\tag{1.3}
\]

All sums in this section are finite.

\begin{center}\rule{0.5\linewidth}{0.5pt}\end{center}

\hypertarget{native-cut-tail-mass}{%
\section{2. Native cut-tail mass}\label{native-cut-tail-mass}}

For a scalar \(z=r+t\iota_\Sigma\), define

\[
D_\Sigma(z)=|r|_\Sigma+|t|_\Sigma.
\tag{2.1}
\]

This is a rational cut gauge. It is not introduced as an absolute value on
\(\mathbb C\).

For a finite path element

\[
a=\sum_\gamma a_\gamma\gamma,
\]

define its cut-tail mass

\[
\boxed{
\mathcal M_\Sigma(a)
=
\sum_\gamma D_\Sigma(a_\gamma).
}
\tag{2.2}
\]

\hypertarget{lemma-2.1-scalar-cut-mass-product-inequality}{%
\subsection{Lemma 2.1 --- Scalar cut-mass product inequality}\label{lemma-2.1-scalar-cut-mass-product-inequality}}

For all \(z,w\in\mathbb K_\Sigma^0\),

\[
D_\Sigma(zw)
\le
D_\Sigma(z)D_\Sigma(w).
\tag{2.3}
\]

\hypertarget{proof}{%
\subsection{Proof}\label{proof}}

Write

\[
z=r+t\iota_\Sigma,
\qquad
w=s+u\iota_\Sigma.
\]

Then

\[
zw=(rs-tu)+(ru+ts)\iota_\Sigma.
\]

The ordered cut-rational triangle inequality gives

\[
\begin{aligned}
D_\Sigma(zw)
&=|rs-tu|_\Sigma+|ru+ts|_\Sigma\\
&\le |r||s|+|t||u|+|r||u|+|t||s|\\
&=(|r|+|t|)(|s|+|u|).
\end{aligned}
\]

QED.

\hypertarget{theorem-2.2-native-cut-tail-convolution-inequality}{%
\subsection{Theorem 2.2 --- Native cut-tail convolution inequality}\label{theorem-2.2-native-cut-tail-convolution-inequality}}

For all finite path elements,

\[
\boxed{
\mathcal M_\Sigma(a\star b)
\le
\mathcal M_\Sigma(a)\mathcal M_\Sigma(b).
}
\tag{2.4}
\]

\hypertarget{proof-1}{%
\subsection{Proof}\label{proof-1}}

Using \texttt{(2.3)} and the finite triangle inequality,

\[
\begin{aligned}
\mathcal M_\Sigma(a\star b)
&=
\sum_\gamma
D_\Sigma\left(
\sum_{\alpha\beta=\gamma}a_\alpha b_\beta
\right)\\
&\le
\sum_{\alpha,\beta\;\mathrm{composable}}
D_\Sigma(a_\alpha)D_\Sigma(b_\beta)\\
&\le
\left(\sum_\alpha D_\Sigma(a_\alpha)\right)
\left(\sum_\beta D_\Sigma(b_\beta)\right).
\end{aligned}
\]

QED.

Dagger preserves cut-tail mass exactly:

\[
\mathcal M_\Sigma(a^\dagger)=\mathcal M_\Sigma(a).
\tag{2.5}
\]

\begin{center}\rule{0.5\linewidth}{0.5pt}\end{center}

\hypertarget{native-cut-tail-completion}{%
\section{3. Native cut-tail completion}\label{native-cut-tail-completion}}

A sequence \((a_n)\) in \(\mathscr P_\Sigma^0\) is \textbf{cut-tail Cauchy} when for
every positive cut rational \(\varepsilon\), there exists \(N\) such that

\[
\mathcal M_\Sigma(a_m-a_n)<\varepsilon
\qquad(m,n\ge N).
\tag{3.1}
\]

Two cut-tail Cauchy sequences are equivalent when

\[
\mathcal M_\Sigma(a_n-b_n)\longrightarrow0
\tag{3.2}
\]

in rational-threshold language.

Define

\[
\boxed{
\mathfrak A_\Sigma^{\mathrm{tail}}
=
\{\text{cut-tail Cauchy presentations}\}/\sim_{\mathcal M}.
}
\tag{3.3}
\]

\hypertarget{lemma-3.1-cauchy-presentations-are-mass-bounded}{%
\subsection{Lemma 3.1 --- Cauchy presentations are mass bounded}\label{lemma-3.1-cauchy-presentations-are-mass-bounded}}

Every cut-tail Cauchy sequence \((a_n)\) admits a positive rational bound
\(B\) such that

\[
\mathcal M_\Sigma(a_n)\le B
\]

for every \(n\).

\hypertarget{proof-2}{%
\subsection{Proof}\label{proof-2}}

Choose \(N\) such that

\[
\mathcal M_\Sigma(a_n-a_N)<1
\qquad(n\ge N).
\]

Then

\[
\mathcal M_\Sigma(a_n)
\le
\mathcal M_\Sigma(a_N)+1
\]

for \(n\ge N\). Enlarge the bound to include the finitely many earlier terms.
QED.

\hypertarget{theorem-3.2-native-cut-tail-algebra}{%
\subsection{Theorem 3.2 --- Native cut-tail algebra}\label{theorem-3.2-native-cut-tail-algebra}}

Termwise addition, convolution, and dagger descend to
\(\mathfrak A_\Sigma^{\mathrm{tail}}\). The resulting object is a complete
involutive path algebra.

\hypertarget{proof-3}{%
\subsection{Proof}\label{proof-3}}

Let \((a_n)\) and \((b_n)\) be cut-tail Cauchy and mass bounded by \(A\) and
\(B\). Then

\[
\begin{aligned}
\mathcal M_\Sigma(a_n\star b_n-a_m\star b_m)
&\le
\mathcal M_\Sigma((a_n-a_m)\star b_n)\\
&\quad+
\mathcal M_\Sigma(a_m\star(b_n-b_m))\\
&\le
B\mathcal M_\Sigma(a_n-a_m)
+A\mathcal M_\Sigma(b_n-b_m).
\end{aligned}
\tag{3.4}
\]

Thus termwise convolution is Cauchy. The same estimate shows independence of
representatives. Dagger is immediate from \texttt{(2.5)}. Associativity and the star
law hold termwise on the finite core and therefore on equivalence classes.

Completeness follows from the standard Cauchy-presentation diagonal argument,
but every neighborhood and estimate used in that argument is expressed through
\(\mathcal M_\Sigma\) and positive cut-rational thresholds. No pre-existing
\(L^1\) space is used to define the object. QED.

The phrase \textbf{cut-tail completion} refers to \texttt{(3.3)}. Any later identification
with an ordinary sequence space is a separate coordinate theorem.

\begin{center}\rule{0.5\linewidth}{0.5pt}\end{center}

\hypertarget{native-closed-loop-recognition-energy}{%
\section{4. Native closed-loop recognition energy}\label{native-closed-loop-recognition-energy}}

Let \(\Phi_\Sigma\) extract stationary zero-residue paths. Define

\[
\langle a,b\rangle_\Sigma^0
=
\Phi_\Sigma(a^\dagger\star b).
\tag{4.1}
\]

The typed path rules give the exact coefficient identity

\[
\boxed{
\langle a,b\rangle_\Sigma^0
=
\sum_\gamma a_\gamma^\dagger b_\gamma.
}
\tag{4.2}
\]

Define recognition energy

\[
\boxed{
\mathcal E_\Sigma(a)
=
\langle a,a\rangle_\Sigma^0
=
\sum_\gamma N_\Sigma(a_\gamma).
}
\tag{4.3}
\]

It is a nonnegative radial cut rational, and

\[
\mathcal E_\Sigma(a)=0
\Longleftrightarrow
a=0.
\tag{4.4}
\]

\hypertarget{lemma-4.1-native-addition-control}{%
\subsection{Lemma 4.1 --- Native addition control}\label{lemma-4.1-native-addition-control}}

For finite paths,

\[
\boxed{
\mathcal E_\Sigma(a+b)
\le
2\mathcal E_\Sigma(a)+2\mathcal E_\Sigma(b).
}
\tag{4.5}
\]

\hypertarget{proof-4}{%
\subsection{Proof}\label{proof-4}}

For every pair of cut scalars,

\[
2N_\Sigma(x)+2N_\Sigma(y)-N_\Sigma(x+y)
=N_\Sigma(x-y)\ge0.
\tag{4.6}
\]

Sum \texttt{(4.6)} over path coefficients. QED.

\begin{center}\rule{0.5\linewidth}{0.5pt}\end{center}

\hypertarget{native-finite-cauchy-schwarz-identity}{%
\section{5. Native finite Cauchy-Schwarz identity}\label{native-finite-cauchy-schwarz-identity}}

The required recognition estimate is not imported from a Hilbert space.

Let \(x_1,\ldots,x_n\) be cut scalars and let
\(d_1,\ldots,d_n\) be positive cut rationals. Put

\[
D=\sum_i d_i.
\]

\hypertarget{lemma-5.1-weighted-cut-square-identity}{%
\subsection{Lemma 5.1 --- Weighted cut-square identity}\label{lemma-5.1-weighted-cut-square-identity}}

\[
\boxed{
D\sum_i\frac{N_\Sigma(x_i)}{d_i}
-
N_\Sigma\left(\sum_i x_i\right)
=
\sum_{i<j}
\frac{N_\Sigma(d_jx_i-d_ix_j)}{d_id_j}.
}
\tag{5.1}
\]

\hypertarget{proof-5}{%
\subsection{Proof}\label{proof-5}}

Expand both sides using

\[
N_\Sigma(z)=z^\dagger z.
\]

The coefficient of \(N_\Sigma(x_i)\) on each side is

\[
\frac{D-d_i}{d_i},
\]

and every mixed pair contributes

\[
-x_i^\dagger x_j-x_j^\dagger x_i.
\]

Thus the expressions agree exactly. Every term on the right is a positive
radial square. QED.

\hypertarget{corollary-5.2-native-coefficient-pairing-inequality}{%
\subsection{Corollary 5.2 --- Native coefficient pairing inequality}\label{corollary-5.2-native-coefficient-pairing-inequality}}

For finite path elements,

\[
\boxed{
N_\Sigma(\langle a,b\rangle_\Sigma^0)
\le
\mathcal E_\Sigma(a)\mathcal E_\Sigma(b).
}
\tag{5.2}
\]

This also follows from the unweighted finite identity

\[
\mathcal E_\Sigma(a)\mathcal E_\Sigma(b)
-
N_\Sigma(\langle a,b\rangle_\Sigma^0)
=
\sum_{\gamma<\delta}
N_\Sigma(a_\gamma b_\delta-a_\delta b_\gamma).
\tag{5.3}
\]

No square-root norm is required.

\begin{center}\rule{0.5\linewidth}{0.5pt}\end{center}

\hypertarget{native-recognition-cauchy-completion}{%
\section{6. Native recognition-Cauchy completion}\label{native-recognition-cauchy-completion}}

A sequence \((\xi_n)\) of finite paths is \textbf{recognition Cauchy} when for every
positive cut rational \(\varepsilon\), there exists \(N\) such that

\[
\mathcal E_\Sigma(\xi_m-\xi_n)<\varepsilon
\qquad(m,n\ge N).
\tag{6.1}
\]

Two recognition-Cauchy presentations are equivalent when

\[
\mathcal E_\Sigma(\xi_n-\eta_n)\longrightarrow0.
\tag{6.2}
\]

Define

\[
\boxed{
\mathfrak R_\Sigma
=
\{\text{recognition-Cauchy presentations}\}/\sim_{\mathcal E}.
}
\tag{6.3}
\]

\hypertarget{theorem-6.1-native-recognition-completion}{%
\subsection{Theorem 6.1 --- Native recognition completion}\label{theorem-6.1-native-recognition-completion}}

The object \(\mathfrak R_\Sigma\) is a complete module over the completed cut
field. The pairing \texttt{(4.1)} extends uniquely to

\[
\langle\cdot,\cdot\rangle_\Sigma:
\mathfrak R_\Sigma\times\mathfrak R_\Sigma
\longrightarrow
\widehat{\mathbb K}_\Sigma,
\tag{6.4}
\]

and satisfies

\[
\langle\xi,\xi\rangle_\Sigma\ge0,
\qquad
\langle\xi,\xi\rangle_\Sigma=0
\Longleftrightarrow
\xi=0.
\tag{6.5}
\]

\hypertarget{proof-6}{%
\subsection{Proof}\label{proof-6}}

Addition is continuous by \texttt{(4.5)}, and scalar multiplication obeys

\[
\mathcal E_\Sigma(z\xi)
=N_\Sigma(z)\mathcal E_\Sigma(\xi).
\tag{6.6}
\]

Recognition-Cauchy sequences are energy bounded by the same argument as Lemma
3.1, using \texttt{(4.5)}.

For Cauchy presentations \((\xi_n)\), \((\eta_n)\), the difference

\[
\langle\xi_n,\eta_n\rangle^0
-
\langle\xi_m,\eta_m\rangle^0
\]

is the sum of

\[
\langle\xi_n-\xi_m,\eta_n\rangle^0
\quad\text{and}\quad
\langle\xi_m,\eta_n-\eta_m\rangle^0.
\]

Use \texttt{(4.6)} and \texttt{(5.2)}. Energy boundedness makes both terms converge to zero
in cut norm square. Therefore the pairing sequence is cut-Cauchy and independent
of representatives.

Positivity follows from \texttt{(4.3)}. If the completed self-pairing is zero, then
\(\mathcal E_\Sigma(\xi_n)\) converges to zero, so the presentation is
recognition-null and represents the zero class.

Completeness follows from the Cauchy-presentation diagonal construction using
only the entourages \texttt{(6.1)}. QED.

This object is called the \textbf{recognition completion}. The theorem does not use
the word Hilbert as a definition. A later chart may prove that its coordinate
image is an ordinary Hilbert space.

\begin{center}\rule{0.5\linewidth}{0.5pt}\end{center}

\hypertarget{native-tail-action-on-recognition-completion}{%
\section{7. Native tail action on recognition completion}\label{native-tail-action-on-recognition-completion}}

\hypertarget{theorem-7.1-finite-native-action-bound}{%
\subsection{Theorem 7.1 --- Finite native action bound}\label{theorem-7.1-finite-native-action-bound}}

For finite paths \(a\) and \(\xi\),

\[
\boxed{
\mathcal E_\Sigma(a\star\xi)
\le
\mathcal M_\Sigma(a)^2\mathcal E_\Sigma(\xi).
}
\tag{7.1}
\]

\hypertarget{proof-7}{%
\subsection{Proof}\label{proof-7}}

Fix an output arrow \(\gamma\). For each decomposition
\(\alpha\beta=\gamma\), put

\[
x_\alpha=a_\alpha\xi_\beta,
\qquad
d_\alpha=D_\Sigma(a_\alpha).
\]

Terms with \(a_\alpha=0\) are omitted, so \(d_\alpha>0\). Lemma 5.1 gives

\[
N_\Sigma\left(
\sum_{\alpha\beta=\gamma}a_\alpha\xi_\beta
\right)
\le
\mathcal M_\Sigma(a)
\sum_{\alpha\beta=\gamma}
\frac{N_\Sigma(a_\alpha\xi_\beta)}{d_\alpha}.
\tag{7.2}
\]

Since

\[
N_\Sigma(a_\alpha)
\le D_\Sigma(a_\alpha)^2=d_\alpha^2,
\]

we obtain

\[
N_\Sigma((a\star\xi)_\gamma)
\le
\mathcal M_\Sigma(a)
\sum_{\alpha\beta=\gamma}
 d_\alpha N_\Sigma(\xi_\beta).
\tag{7.3}
\]

Sum over \(\gamma\). Every composable pair produces one output arrow, and the
sum over composable pairs is bounded by the unrestricted product of the two
finite positive sums. Hence

\[
\mathcal E_\Sigma(a\star\xi)
\le
\mathcal M_\Sigma(a)^2\mathcal E_\Sigma(\xi).
\]

QED.

\hypertarget{theorem-7.2-completed-native-action}{%
\subsection{Theorem 7.2 --- Completed native action}\label{theorem-7.2-completed-native-action}}

Convolution extends uniquely to an action

\[
\boxed{
\mathfrak A_\Sigma^{\mathrm{tail}}
\times
\mathfrak R_\Sigma
\longrightarrow
\mathfrak R_\Sigma.
}
\tag{7.4}
\]

The completed action obeys the native bound \texttt{(7.1)} in cut-completed radial
coordinates.

\hypertarget{proof-8}{%
\subsection{Proof}\label{proof-8}}

Let \((a_n)\) be mass Cauchy and \((\xi_n)\) recognition Cauchy. By \texttt{(4.5)}
and \texttt{(7.1)},

\[
\begin{aligned}
\mathcal E_\Sigma(a_n\star\xi_n-a_m\star\xi_m)
&\le
2\mathcal E_\Sigma((a_n-a_m)\star\xi_n)\\
&\quad+
2\mathcal E_\Sigma(a_m\star(\xi_n-\xi_m))\\
&\le
2\mathcal M_\Sigma(a_n-a_m)^2
  \mathcal E_\Sigma(\xi_n)\\
&\quad+
2\mathcal M_\Sigma(a_m)^2
  \mathcal E_\Sigma(\xi_n-\xi_m).
\end{aligned}
\tag{7.5}
\]

Mass and energy boundedness show that \texttt{(7.5)} tends to zero. The same estimate
proves independence of representatives. QED.

This is native boundedness: preservation of recognition-Cauchy neighborhoods by
a cut-tail mass estimate. No operator norm is primitive.

\begin{center}\rule{0.5\linewidth}{0.5pt}\end{center}

\hypertarget{residue-is-not-lost-by-completion}{%
\section{8. Residue is not lost by completion}\label{residue-is-not-lost-by-completion}}

Let \(r\ne s\) be residue channels over the same endpoints. Then

\[
w=\delta_{(x,y,r)}-\delta_{(x,y,s)}
\]

has

\[
\mathcal E_\Sigma(w)=2,
\tag{8.1}
\]

while the endpoint-only shadow sends \(w\) to zero.

Thus neither the cut-tail completion nor the recognition completion erases seam
memory. Blindness enters only through a later quotient map.

\begin{center}\rule{0.5\linewidth}{0.5pt}\end{center}

\hypertarget{explicit-infinite-native-presentation}{%
\section{9. Explicit infinite native presentation}\label{explicit-infinite-native-presentation}}

For the integer object atlas with two residue channels, define

\[
s_N
=
\sum_{n=0}^{N}
2^{-(n+1)}
\delta_{(0,n,n\bmod2)}.
\tag{9.1}
\]

For \(M>N\),

\[
\mathcal M_\Sigma(s_M-s_N)
=
\sum_{n=N+1}^{M}2^{-(n+1)}
<2^{-(N+1)},
\tag{9.2}
\]

and

\[
\mathcal E_\Sigma(s_M-s_N)
=
\sum_{n=N+1}^{M}4^{-(n+1)}
<\frac{1}{3\,4^{N+1}}.
\tag{9.3}
\]

Therefore the same finite presentation converges in both native completion
relations and retains alternating residue over infinitely many targets.

\begin{center}\rule{0.5\linewidth}{0.5pt}\end{center}

\hypertarget{what-has-been-removed}{%
\section{10. What has been removed}\label{what-has-been-removed}}

The construction order is now

\[
\boxed{
\text{cut scalars}
\to
\text{finite typed paths}
\to
\mathcal M_\Sigma,\mathcal E_\Sigma
\to
\text{two native Cauchy quotients}
\to
\text{native completed action}.
}
\]

The following are not used to define or prove this chain:

\begin{verbatim}
ordinary complex numbers;
ell1 or ell2 spaces;
Banach or Hilbert completeness as imported objects;
operator norms;
Haar measure;
Minkowski inequality;
Young inequality;
spectral theory.
\end{verbatim}

Finite algebraic cut-square identities replace the imported analytic bounds.

\begin{center}\rule{0.5\linewidth}{0.5pt}\end{center}

\hypertarget{coordinate-shadow-boundary}{%
\section{11. Coordinate-shadow boundary}\label{coordinate-shadow-boundary}}

This theorem does not deny that a later coordinate map may identify

\begin{verbatim}
native cut-tail completion       with an ordinary absolute-summability space;
native recognition completion   with an ordinary square-summability space;
native completed action          with a bounded regular representation.
\end{verbatim}

Those identifications must be proved \textbf{after} the native objects exist. They
are coordinate-shadow theorems, not definitions.

\begin{center}\rule{0.5\linewidth}{0.5pt}\end{center}

\hypertarget{exact-certificate}{%
\section{12. Exact certificate}\label{exact-certificate}}

The exact dense-core certificate has expected status

\begin{verbatim}
PASS_NATIVE_CUT_TAIL_AND_RECOGNITION_CAUCHY_COMPLETION_CORE
\end{verbatim}

It verifies:

\begin{verbatim}
quarter-turn scalar algebra;
dagger and radial squares;
cut-tail scalar and convolution inequalities;
finite path associativity and dagger reversal;
closed-loop coefficient identity;
native addition control;
weighted cut Cauchy-Schwarz identity;
native action-energy bound;
explicit mass and energy Cauchy tails;
nonzero residue with zero endpoint shadow;
absence of classical-analysis imports in the foundational module.
\end{verbatim}

The certificate is an implementation witness. The proof above establishes the
abstract completion theorem.

\begin{center}\rule{0.5\linewidth}{0.5pt}\end{center}

\hypertarget{remaining-gates}{%
\section{13. Remaining gates}\label{remaining-gates}}

\begin{verbatim}
T01-A NATIVE CUT-TAIL COMPLETION                   PROVED
T01-B NATIVE RECOGNITION-CAUCHY COMPLETION        PROVED
T01-C NATIVE COMPLETED PATH ACTION                PROVED

ACTUAL MP PATH-ATLAS IDENTIFICATION               OPEN
ACTUAL MP REFINEMENT CONTINUITY                   OPEN
NATIVE CONTINUOUS REFINEMENT MEASURE              OPEN
COORDINATE SHADOW TO ORDINARY ell1/ell2           OPEN
NATIVE STATE-PACKET GENERATION                    OPEN
TARGET FAITHFULNESS                               OPEN
NATIVE SHUNYA                                     OPEN
RIEMANN HYPOTHESIS                                OPEN
\end{verbatim}

The next native reconstruction gate is a refinement-count measure derived from
path partitions before any Haar measure is named.

\begin{flushright}\small\texttt{RH-Framework/theorems/T01\_NATIVE\_SUMMABILITY\_AND\_RECOGNITION\_COMPLETION.md}\end{flushright}
